% =============================================================================
% Szkic artykułu wygenerowany przez analiza_odpowiedzi_naukowa.py.
% Plik jest NADPISYWANY przy każdym uruchomieniu skryptu -- zmiany wprowadzaj w kopii.
% Liczby: macros.tex; tabele: \TabDir; ryciny: \FigDir; literatura: references.bib.
% Miejsca do uzupełnienia: \fillin{...}. Zdania interpretujące kierunek wyniku oznaczono
% komentarzem "% SPRAWDŹ" -- zweryfikuj je po zmianie danych lub parametrów.
% Kompilacja podglądu: pdflatex main && bibtex main && pdflatex main && pdflatex main
% =============================================================================

\begin{abstract}
\noindent\textbf{Background.} Spaced retrieval practice improves retention in laboratory and
educational settings, but evidence from long-running workplace programmes is scarce.
\textbf{Methods.} We analysed \NResponses{} responses given by \NEmployees{} employees to \NItems{}
multiple-choice items in \NModules{} modules of a quiz-based microlearning programme
(\DateStart--\DateEnd). Each item was re-presented to the same employee at a fixed, non-adaptive
interval (median \GapMedian{} days). Correctness was modelled with logistic regression including
employee and item fixed effects; descriptive estimates carry cluster-bootstrap confidence intervals.
\textbf{Results.} Accuracy rose from \AccAttOne\% at the first attempt to \AccAttMax\% at the
\AttMax{}th and later attempts. Each doubling of the number of attempts multiplied the odds of a
correct answer by \ORMOneAttempt{} (95\% CI \ORMOneAttemptCI; range across ten sensitivity analyses
\ORRange). Longer intervals since the previous attempt were associated with lower accuracy
(OR \ORMTwoGap{} per doubling of the interval). Repeated errors concentrated on the same distractor
(\SameDistractorPct\% vs \SameDistractorExpectedPop\% expected). Accuracy on newly encountered items
did not improve with tenure (OR \ORMThreeTenure{} per year). % SPRAWDŹ
\textbf{Conclusions.} Repeated retrieval in routine workplace training was accompanied by steady,
item-specific gains and measurable forgetting between repetitions. Adaptive scheduling, targeted
feedback on persistent misconceptions and revision of poorly functioning items could increase the
programme's effectiveness.

\medskip\noindent\textbf{Keywords:} spaced repetition; retrieval practice; microlearning;
workplace training; learning analytics; item analysis
\end{abstract}

% -----------------------------------------------------------------------------
\section{Introduction}

Retrieving information from memory strengthens its later retention more than restudying it
(the testing effect), and practice distributed over time is retained better than massed practice
(the spacing effect) \citep{roediger2006,cepeda2006,cepeda2008}. Spaced education, in which short
multiple-choice questions are delivered repeatedly over weeks or months, improved knowledge
acquisition and retention in randomised trials among medical trainees
\citep{kerfoot2007urology,kerfoot2007mededuc,kerfoot2009}, and systematic reviews report small to
moderate benefits for health professionals \citep{phillips2019,martinengo2024}.

In corporate settings, short quiz-based \emph{microlearning} formats are widely used, but the
evidence base is dominated by small or grey-literature studies \citep{taylor2022}. Evaluations of
mandatory organisational training, for example in security awareness, frequently report modest or
decaying effects \citep{reinheimer2020,lain2022,ho2025}. At the same time, the log data generated by
such platforms allow learning to be observed at the level of single responses, as in learner models
developed for intelligent tutoring and language-learning systems
\citep{corbett1994,pavlik2009,settles2016,tabibian2019,koedinger2023}.

\fillin{One or two sentences on the organisation and the training content.} We used three years of
response logs from a workplace microlearning programme to address the following questions:
(RQ1) does accuracy increase with repeated retrieval of the same item by the same employee;
(RQ2) how is accuracy related to the outcome of, and the time since, the previous attempt;
(RQ3) are repeated errors persistent, indicating misconceptions rather than guessing;
(RQ4) does improvement extend to newly encountered items; and
(RQ5) how well do the items function psychometrically?

% -----------------------------------------------------------------------------
\section{Methods}

\subsection{Setting and training programme}
\fillin{Organisation, sector, target group and training topic; whether participation was mandatory;
delivery channel (e-mail, app, intranet); whether the correct answer and an explanation were shown
after each response; how items were authored and reviewed.}
Employees were presented with short multiple-choice questions (median \OptionsMedian{} response
options per item) grouped into thematic modules. In the observed data employees answered a median of
\DailyMedian{} questions per active day, \WeekdayPct\% of responses were given on weekdays, and each
item was presented again to the same employee after a median of \GapMedian{} days
(IQR \GapQOne--\GapQThree).

\subsection{Data}
The platform log contained, for each response, a record identifier, the module, the item, the option
chosen, a pseudonymised employee identifier, a time stamp and the correctness of the response.
Records with missing or invalid values were removed (\NRejected{} records). Test module(s)
(\ExcludedModules; \NExcluded{} responses) were excluded. The analytic dataset comprised
\NResponses{} responses (Table~\ref{tab:dataset}). \fillin{Ethics approval or waiver; legal basis for
processing employee data (GDPR); how pseudonymisation was performed.}

\subsection{Definitions}
A \emph{sequence} is the chronologically ordered series of attempts of one employee at one item;
the \emph{attempt number} is the position of a response within its sequence (ties in time were
broken by record identifier). \emph{Accuracy} is the proportion of correct responses. The
\emph{interval} is the time in days since the previous attempt in the same sequence. \emph{Tenure} is
the time since the employee's first response in the programme. Employee-level change was defined as
the difference in the proportion of correct answers between attempt~3 and attempt~1, computed across
the employee's own sequences with at least three attempts.

\subsection{Statistical analysis}
Descriptive accuracy is reported with 95\% confidence intervals from a cluster bootstrap that
resampled employees (\BootReps{} replicates), which accounts for the dependence of responses within
persons. Because learners who continue to later attempts may differ from those who stop, aggregated
learning curves can be distorted by attrition \citep{murray2013,goutte2018}; learning curves were
therefore additionally computed for a balanced cohort of sequences with at least \CohortK{}
attempts. The relationship between error rate and attempt number was summarised by a power function
\citep{newell1981}. Because the scheduler was not adaptive, we checked whether re-presentation
depended on the previous outcome.

Correctness was modelled with logistic regression including fixed effects for employee and item and
with standard errors clustered by employee, so that the attempt effect is estimated within persons
and within items. Model~1 included the attempt number on the log$_2$ scale, so the odds ratio (OR)
refers to a doubling of the number of attempts. Model~2, fitted to repeated attempts, added the
outcome of the previous attempt, the interval since the previous attempt (log$_2$ days) and tenure,
analogous to performance-factors and half-life regression models of learner logs
\citep{pavlik2009,settles2016}. Model~3, fitted to first attempts only, tested whether accuracy on
newly encountered items changed with tenure (RQ4). Sensitivity analyses used a balanced cohort,
truncation at 10 attempts, strata of sequence length (addressing pattern-mixture bias), exclusion of
low-activity employees, adjustment for tenure, generalised estimating equations with an
exchangeable working correlation \citep{liang1986}, and a mixed model with crossed random
intercepts for employees and items \citep{baayen2008}, from which latent-scale intraclass
correlations (ICC) and marginal and conditional $R^2$ were derived \citep{nakagawa2013}.

Transitions between consecutive attempts were tabulated. For consecutive errors, the proportion of
repeated choices of the same distractor was compared with the proportion expected under random
choice among the observed distractors and under choice proportional to each distractor's popularity
for the item. Employee-level change was tested with the sign test and the Wilcoxon signed-rank test
(rank-biserial correlation as effect size). Items were described with classical test theory indices
computed on first attempts: difficulty (proportion correct), discrimination (item--rest correlation
with the employee's accuracy on the remaining items) and distractor functioning, with distractors
chosen by fewer than 5\% of respondents classified as non-functioning \citep{tarrant2009,gierl2017}.
Per-item trends across attempts were tested with logistic regression with employee-clustered
standard errors, controlling the false discovery rate at \Alpha{} \citep{benjamini1995}.
Analyses were performed in Python \VerPython{} with pandas \VerPandas, NumPy \VerNumpy,
SciPy \VerScipy, statsmodels \VerStatsmodels{} and Matplotlib \VerMatplotlib{} (random seed \Seed).

% -----------------------------------------------------------------------------
\section{Results}

\subsection{Sample and engagement}
The dataset comprised \NResponses{} responses from \NEmployees{} employees to \NItems{} items in
\NModules{} modules, collected between \DateStart{} and \DateEnd{} (Table~\ref{tab:dataset}).
Employees gave a median of \RespPerEmpMedian{} responses (IQR \RespPerEmpQOne--\RespPerEmpQThree)
to \ItemsPerEmpMedian{} different items on \ActiveDaysMedian{} active days, and remained in the
programme for a median of \SpanMonthsMedian{} months; \ActiveLastQuarterPct\% were active in the
last 90 days of observation. Activity varied over time, with periods without responses
(Figure~\ref{fig:activity}). \fillin{Explain the gaps, e.g. programme pauses or module launches.}
Overall accuracy was \AccOverall\% and first-attempt accuracy \AccFirst\%.

The re-presentation schedule did not depend on the previous outcome: the probability of a further
attempt was \ReaskAfterCorrect\% after a correct and \ReaskAfterIncorrect\% after an incorrect
response, and the median interval to the next attempt was \GapAfterCorrect{} and
\GapAfterIncorrect{} days, respectively. Continuation of sequences was therefore not selected on
performance. % SPRAWDŹ

\input{\TabDir dataset}

\begin{figure*}[tbp]
\centering
\includegraphics{\FigDir fig_activity}
\caption{Monthly number of responses (top) and of active employees (bottom) over the observation
period.}
\label{fig:activity}
\end{figure*}

\subsection{Learning across repeated attempts (RQ1)}
Accuracy increased steadily with the attempt number, from \AccAttOne\% \AccAttOneCI{} at the
first attempt to \AccAttMax\% \AccAttMaxCI{} at the \AttMax{}th and later attempts
(Figure~\ref{fig:learning_curve}, Table~\ref{tab:learning_curve}). The increase was of the same size
in the balanced cohort of \CohortN{} sequences with at least \CohortK{} attempts (\CohortAccOne\%
\CohortAccOneCI{} at attempt~1 vs \CohortAccK\% \CohortAccKCI{} at attempt~\CohortK), so it is not
explained by selective continuation of better-performing employees. The error rate followed a power
function of the attempt number (exponent \PowerLawB, $R^2$ = \PowerLawRsq), consistent with the power
law of practice \citep{newell1981}.

In Model~1, each doubling of the number of attempts was associated with \ORMOneAttempt-fold higher
odds of a correct answer (95\% CI \ORMOneAttemptCI; Table~\ref{tab:regression}). Items accounted for a
larger share of the variance than employees (latent ICC \ICCItem{} vs \ICCEmployee; marginal
$R^2$ = \RsqMarginal, conditional $R^2$ = \RsqConditional).

\begin{figure}[tbp]
\centering
\includegraphics{\FigDir fig_learning_curve}
\caption{Accuracy by attempt number within employee--item sequences. Circles: all sequences
available at each attempt (their number shrinks as attempts increase). Squares: balanced cohort,
i.e.\ the same \CohortN{} sequences that reached at least \CohortK{} attempts, followed over
attempts 1--\CohortK{} only, so that every point is based on identical sequences. Error bars: 95\% CI from a
cluster bootstrap over employees (\BootReps{} replicates). Dashed line: power-law fit to the error
rate. Attempts \MaxAttemptPooled{} and higher are pooled.}
\label{fig:learning_curve}
\end{figure}

\input{\TabDir learning_curve}

\input{\TabDir regression}

\subsection{Previous outcome, interval and forgetting (RQ2)}
Controlling for the attempt number, a correct previous response was associated with higher odds of a
correct answer (Model~2: OR \ORMTwoPrev, 95\% CI \ORMTwoPrevCI), and longer intervals since the
previous attempt with lower odds (OR \ORMTwoGap{} per doubling of the interval, 95\% CI
\ORMTwoGapCI; Table~\ref{tab:regression}). Among previously correct responses, accuracy decreased from
\RetShortCorrect\% when the item was repeated within 14 days to \RetLongCorrect\% after 240 days or
more; previously incorrect responses were corrected in \RetShortIncorrect\% and \RetLongIncorrect\% of
cases, respectively (Figure~\ref{fig:retention}, Table~\ref{tab:retention_interval}). This pattern is
consistent with forgetting between repetitions \citep{cepeda2006,cepeda2008}. % SPRAWDŹ

\begin{figure}[tbp]
\centering
\includegraphics{\FigDir fig_retention_interval}
\caption{Accuracy by the interval since the previous attempt at the same item, separately for
previously correct (circles) and previously incorrect (squares) responses. Error bars: 95\% CI from a
cluster bootstrap over employees.}
\label{fig:retention}
\end{figure}

\input{\TabDir retention_interval}

\subsection{Transitions and persistent errors (RQ3)}
Of previously correct responses, \RetainPct\% \RetainCI{} were correct again at the next attempt
(\ForgetPct\% were forgotten); of previously incorrect responses, \RecoverPct\% \RecoverCI{} were
corrected and \PersistPct\% remained incorrect (Table~\ref{tab:transitions}). When an error was
repeated (\SameDistractorN{} cases), employees chose the same distractor in \SameDistractorPct\%
\SameDistractorCI{} of cases, compared with \SameDistractorExpected\% expected under random choice
and \SameDistractorExpectedPop\% expected given the popularity of each distractor, indicating
stable misconceptions rather than guessing. % SPRAWDŹ

\input{\TabDir transitions}

\subsection{Employee-level change}
Among the \EmpN{} employees (\EmpSharePct\% of all) with at least one sequence of three or more
attempts, \EmpImproved{} improved, \EmpWorsened{} worsened and \EmpUnchanged{} did not change between
the first and third attempt (sign test $p$ \SignP; Wilcoxon $p$ \WilcoxonP; rank-biserial
$r$ = \RankBiserial). The mean change was \EmpMeanChange{} percentage points (95\% CI
\EmpMeanChangeCI; Figure~\ref{fig:employee_change}, Table~\ref{tab:employee_change}).

\begin{figure}[tbp]
\centering
\includegraphics{\FigDir fig_employee_change}
\caption{Distribution of the employee-level change in accuracy between attempt~3 and attempt~1
across the employee's own sequences with at least three attempts ($n$ = \EmpN{} employees; bins of
5 percentage points centred on zero).}
\label{fig:employee_change}
\end{figure}

\input{\TabDir employee_change}

\subsection{Transfer to new items (RQ4)}
Accuracy at the first encounter with a new item did not increase with tenure (Model~3: OR
\ORMThreeTenure{} per year, 95\% CI \ORMThreeTenureCI; Table~\ref{tab:regression}), and tenure was not
positively associated with accuracy on repeated attempts once the attempt number was controlled
(Model~2: OR \ORMTwoTenure, 95\% CI \ORMTwoTenureCI). The gains therefore appear to be specific to
the repeated items. % SPRAWDŹ

\subsection{Item analysis (RQ5)}
The median item difficulty was \ItemDifficultyMedian\% correct at the
first attempt: \ItemsHard{} items were answered correctly by fewer than 50\%, \ItemsModerate{} by
50--90\% and \ItemsEasy{} by more than 90\% of employees (Table~\ref{tab:items}). The median
item--rest correlation was \ItemDiscMedian; \ItemsLowDisc{} of \ItemsDiscTested{} items had
correlations below 0.20 and \ItemsGoodDisc{} of 0.30 or more (Figure~\ref{fig:items}). Errors
concentrated on one distractor (median \TopDistractorMedian\% of an item's errors), and
\NNonFunctional{} of \NDistractors{} observed distractors (\NonFunctionalPct\%) were
non-functioning. Accuracy improved significantly with repetition for \ItemsImproving{} of
\ItemsTestable{} testable items and deteriorated for \ItemsWorsening{} (FDR-adjusted). The most
difficult items are listed in Table~\ref{tab:hardest_items}.

\begin{figure}[tbp]
\centering
\includegraphics{\FigDir fig_item_analysis}
\caption{Item difficulty (accuracy at first attempt) versus discrimination (item--rest
correlation). Triangles: items whose accuracy increased significantly with repetition
(Benjamini--Hochberg $q<\Alpha$). Dotted lines: $r$ = 0.20 and 90\% correct.}
\label{fig:items}
\end{figure}

\input{\TabDir items}

\input{\TabDir hardest_items}

\subsection{Sensitivity analyses}
The attempt effect was stable across specifications, with ORs from \ORRange{}
(Figure~\ref{fig:sensitivity}, Table~\ref{tab:sensitivity}): \ORSensCohort{} in the balanced cohort,
\ORSensLenThree, \ORSensLenFive{} and \ORSensLenEight{} in sequences of 3--4, 5--7 and at least 8
attempts, \ORSensGee{} in the population-averaged GEE model and \ORSensMixed{} in the mixed model.
The estimate was larger after adjustment for tenure (\ORSensTenure), which is correlated with the
attempt number.

\begin{figure}[tbp]
\centering
\includegraphics{\FigDir fig_sensitivity}
\caption{Odds ratio (95\% CI) for a correct answer per doubling of the attempt number across model
specifications. Diamond: main model; vertical line: no effect.}
\label{fig:sensitivity}
\end{figure}

\input{\TabDir sensitivity}

% -----------------------------------------------------------------------------
\section{Discussion}

\subsection{Principal findings}
In three years of routine use of a workplace microlearning programme, accuracy on repeatedly
presented items rose steadily with the number of retrievals, following a power function, and the
effect was robust to attrition, model specification and clustering. Accuracy declined with the
interval since the previous retrieval, previous errors were only partly corrected at the next
attempt, and repeated errors mostly reproduced the same distractor. Accuracy on new items did not
increase with tenure. % SPRAWDŹ

\subsection{Comparison with previous research}
The gradual, regular gains resemble learning curves observed across educational datasets
\citep{koedinger2023} and are consistent with the testing and spacing effects
\citep{roediger2006,cepeda2006}. Randomised trials of spaced education found improved retention
\citep{kerfoot2007mededuc,kerfoot2009}; our observational data extend these findings to a
non-medical workplace setting with long inter-repetition intervals. The negative association between
interval and accuracy agrees with forgetting-curve models fitted to learner logs
\citep{settles2016,tabibian2019}. As in studies that re-used the same items repeatedly
\citep{donker2022}, improvement on repeated items cannot be separated from memory for the specific
question and its answer; the absence of a tenure effect on new items suggests limited transfer.

\subsection{Implications for practice}
The programme re-presented items at the same interval regardless of the previous outcome. Adaptive
schedules that shorten the interval after errors and lengthen it after correct answers use learner
history more efficiently \citep{lindsey2014,settles2016,tabibian2019}. The persistence of specific
distractors points to misconceptions that may call for explanatory feedback rather than repetition
alone. Item analysis identified candidates for revision: items with low discrimination, items
answered correctly by almost everyone at the first attempt (little room for learning) and
non-functioning distractors \citep{tarrant2009,gierl2017}.

\subsection{Limitations}
The study is observational and has no control group, so the gains cannot be attributed causally to
the programme; improvement on repeated items may partly reflect recognition of the item rather than
understanding. Participation and activity varied between employees, and results for sequences with
many attempts describe the more active participants, although the balanced-cohort and
sequence-length analyses did not suggest attrition bias. The high first-attempt accuracy
(\AccFirst\%) limits the room for improvement (ceiling effect). Distractors never chosen are not
visible in the log, so distractor functioning may be overestimated. No information on employees'
roles, prior knowledge or demographics was available. Time stamps carried no time-zone information.
The mixed model included random intercepts only; employee-specific learning rates were not
modelled. The data come from a single organisation. \fillin{Add context-specific limitations.}

% -----------------------------------------------------------------------------
\section{Conclusions}
Routine repeated retrieval in workplace training was accompanied by steady, item-specific
improvement and by forgetting between repetitions. Log data from such programmes can be used to
monitor learning, evaluate items and inform the design of adaptive repetition schedules.

% -----------------------------------------------------------------------------
\section*{Declarations}
\noindent\textbf{Ethics.} \fillin{Approval / waiver and data-protection basis.}\\
\textbf{Data availability.} \fillin{Availability of pseudonymised data and analysis code.}\\
\textbf{Funding.} \fillin{Funding sources.}\\
\textbf{Competing interests.} \fillin{Declaration.}\\
\textbf{Use of AI tools.} \fillin{Declaration of the use of generative AI, if required by the journal.}
